% !TeX program = xelatex
\documentclass[12pt,a4paper]{article}
\usepackage{fontspec}
\setmainfont{Times New Roman}
\usepackage[english,russian]{babel}
\usepackage{amsmath,amssymb,mathrsfs,cancel}
\usepackage[margin=22mm]{geometry}
\usepackage{graphicx,float}
\usepackage[unicode,hidelinks]{hyperref}
\DeclareMathOperator{\grad}{grad}
\DeclareMathOperator{\Div}{div}
\newcommand{\Rey}{\mathrm{Re}}
\setlength{\emergencystretch}{3em}
\sloppy
\interfootnotelinepenalty=10000
\newcommand{\topic}[1]{\par\medskip\noindent\textbf{#1}\par\nobreak\smallskip}
\renewcommand{\thefootnote}{\arabic{footnote})}
\begin{document}
\section*{Вопрос 9}
Вытеснение несмешивающихся жидкостей. Одномерная задача о поршневом вытеснении несмешивающихся жидкостей (модель Лейбензона — Маскета). Устойчивость фронта вытеснения (элементарная теория).

\medskip
\noindent\textbf{Источник:} К. С. Басниев, А. М. Власов, И. Н. Кочина, В. М. Максимов. \textit{Подземная гидравлика}. Москва: Недра, 1986. Глава 8, §§ 1, 2, 4; с. 188--193, 196--198.
\par\bigskip
\begin{quote}\sloppy
Согласно условию неразрывности потока массы элементарные расходы обеих несжимаемых жидкостей через элемент границы раздела, включающий точку $M$, должны быть равны между собой. Отсюда следует, что нормальные составляющие скоростей фильтрации обеих жидкостей будут равны, т. е. $w_{1n}=w_{2n}$. Давление в пласте в точке $M$ также должно быть одинаково для обеих жидкостей, так как при малых скоростях (ниже звуковых) разрыва давления в сплошном потоке быть не может.
\par\nobreak\smallskip
{\small К. С. Басниев, А. М. Власов, И. Н. Кочина, В. М. Максимов. \textit{Подземная гидравлика}. 1986. Глава 8, § 1, с. 188.}\par
\end{quote}

\begin{quote}\sloppy
Линии тока не будут преломляться только в двух случаях --- при прямолинейно-параллельном и плоскорадиальном движениях границы раздела, когда $w_{1\tau}=w_{2\tau}=0$. Эти задачи прежде всего и будут рассмотрены в данной главе. При этом жидкости (нефть и вода) считаются несжимаемыми, взаимно нерастворимыми и химически не реагирующими одна с другой и с пористой средой. Вытеснение нефти водой предполагается происходящим полностью --- так называемое поршневое вытеснение. Задачи, в которых учитывается неполнота вытеснения, рассматриваются в гл. 9.
\par\nobreak\smallskip
{\small К. С. Басниев, А. М. Власов, И. Н. Кочина, В. М. Максимов. \textit{Подземная гидравлика}. 1986. Глава 8, § 1, с. 189.}\par
\end{quote}

\begin{quote}\sloppy
При поршневом вытеснении нефти водой в пористой среде плотности нефти и воды будем считать одинаковыми. Это позволит рассматривать плоскость контакта нефти и воды вертикальной. А различие в вязкостях нефти $\mu_\text{н}$ и воды $\mu_\text{в}$ будем учитывать.

В случае прямолинейно-параллельного движения схема вытеснения представлена на рис. 8.2. На контуре питания и на галерее поддерживаются соответственно постоянные давления $p_\text{к}$ и $p_\text{г}$. Начальное положение контура нефтеносности $x_0$ параллельно галерее и контуру питания.

Обозначим через $x_\text{в}$ текущее расстояние до контура нефтеносности в момент времени $t$ после начала вытеснения, через $L_\text{к}$ --- расстояние от контура питания до галереи, через $p_\text{в},p_\text{н}$ --- давление в любой точке водоносной и нефтеносной части пласта соответственно, через $p(t)$ --- давление на границе раздела вода---нефть, отстоящей от контура питания на расстоянии $x_\text{в}$.
\par\nobreak\smallskip
{\small К. С. Басниев, А. М. Власов, И. Н. Кочина, В. М. Максимов. \textit{Подземная гидравлика}. 1986. Глава 8, § 2, с. 189.}\par
\end{quote}

\begin{quote}\sloppy
При этом изобарами являются линии, параллельные галерее, и каждую изобару можно рассматривать как контур питания или как галерею. На основании формул (8.3) и (8.4) распределение давления и скорость фильтрации в водоносной области можно записать в виде
\begin{equation}p_\text{в}=p_\text{к}-\frac{p_\text{к}-p(t)}{x_\text{в}(t)}x,\quad0\le x\le x_\text{в}(t);\tag{8.5}\label{q9:eq-8-5}\end{equation}
\begin{equation}w_\text{в}=\frac{k}{\mu_\text{в}}\frac{p_\text{к}-p(t)}{x_\text{в}(t)}.\tag{8.6}\label{q9:eq-8-6}\end{equation}
Принимая за контур питания изобару, совпадающую с границей раздела жидкостей, распределение давления и скорость фильтрации в нефтеносной области можно записать следующим образом:
\begin{equation}p_\text{н}=p_\text{г}+\frac{p(t)-p_\text{г}}{L_\text{к}-x_\text{в}(t)}(L_\text{к}-x),\quad x_\text{в}(t)\le x\le L_\text{к};\tag{8.7}\label{q9:eq-8-7}\end{equation}
\begin{equation}w_\text{н}=\frac{k}{\mu_\text{н}}\frac{p(t)-p_\text{г}}{L_\text{к}-x_\text{в}(t)}.\tag{8.8}\label{q9:eq-8-8}\end{equation}
Найдем давление $p(t)$ на границе раздела. Вследствие несжимаемости жидкостей и неразрывности потока линии тока будут иметь вид прямых, параллельных оси $Ox$ (на границе раздела преломления не будет), а скорость фильтрации во всех точках пласта будет одинаковой, т. е. $w_\text{в}=w_\text{н}$.

Тогда из уравнений (8.6) и (8.8) получим
\begin{equation}\frac{k(p_\text{к}-p(t))}{\mu_\text{в}x_\text{в}(t)}=\frac{k(p(t)-p_\text{г})}{\mu_\text{н}(L_\text{к}-x_\text{в}(t))},\tag{8.9}\label{q9:eq-8-9}\end{equation}
откуда давление на границе раздела жидкостей будет
\begin{equation}p(t)=\frac{p_\text{к}\mu_\text{н}(L_\text{к}-x_\text{в}(t))+p_\text{г}\mu_\text{в}x_\text{в}(t)}{\mu_\text{в}x_\text{в}(t)+\mu_\text{н}(L_\text{к}-x_\text{в}(t))}.\tag{8.10}\label{q9:eq-8-10}\end{equation}
\par\nobreak\smallskip
{\small К. С. Басниев, А. М. Власов, И. Н. Кочина, В. М. Максимов. \textit{Подземная гидравлика}. 1986. Глава 8, § 2, с. 190.}\par
\end{quote}

\begin{quote}\sloppy
2. Скорость фильтрации. Подставим (8.10) в (8.6) и (8.8):
\begin{equation}w_\text{в}=w_\text{н}=\frac{k(p_\text{к}-p_\text{г})}{\mu_\text{в}x_\text{в}(t)+\mu_\text{н}(L_\text{к}-x_\text{в}(t))}.\tag{8.13}\label{q9:eq-8-13}\end{equation}
\par\nobreak\smallskip
{\small К. С. Басниев, А. М. Власов, И. Н. Кочина, В. М. Максимов. \textit{Подземная гидравлика}. 1986. Глава 8, § 2, с. 190.}\par
\end{quote}

\begin{quote}\sloppy
5. Закон движения границы раздела $x_\text{в}=x_\text{в}(t)$ находим из соотношения скорости фильтрации и средней скорости движения:
\[w=mv=m\frac{dx_\text{в}}{dt},\]
откуда
\begin{equation}dt=\frac{m}{w}dx_\text{в}=\frac{m}{\dfrac{k(p_\text{к}-p_\text{г})}{\mu_\text{в}x_\text{в}(t)+\mu_\text{н}(L_\text{к}-x_\text{в}(t))}}dx_\text{в}.\tag{8.17}\label{q9:eq-8-17}\end{equation}
Проинтегрировав (8.17) в пределах от 0 до $t$ и от $x_0$ до $x_\text{в}$, получим
\begin{equation}t=\frac{m}{k(p_\text{к}-p_\text{г})}\left[\mu_\text{н}L_\text{к}(x_\text{в}-x_0)-\frac12(\mu_\text{н}-\mu_\text{в})(x_\text{в}^2-x_0^2)\right].\tag{8.18}\label{q9:eq-8-18}\end{equation}
Чтобы найти время полного вытеснения нефти, нужно в формуле (8.18) положить $x_\text{в}=L_\text{к}$. Тогда получим
\[T=\frac{m}{2k(p_\text{к}-p_\text{г})}\left[\mu_\text{в}(L_\text{к}^2-x_0^2)+\mu_\text{н}(L_\text{к}-x_0)^2\right].\]
\par\nobreak\smallskip
{\small К. С. Басниев, А. М. Власов, И. Н. Кочина, В. М. Максимов. \textit{Подземная гидравлика}. 1986. Глава 8, § 2, с. 191.}\par
\end{quote}

\begin{quote}\sloppy
Пусть нефтяная залежь в наклонном пласте (рис. 8.6) имеет горизонтальное начальное положение водонефтяного контакта $A_0B_0$.

При отборе нефти граница раздела вода---нефть будет перемещаться, занимая последовательно положения $A_1B_1$, $A_2B_2$ и т. д. Рассмотрим вопрос об устойчивости движения границы раздела. Скорости фильтрации каждой жидкости согласно закону Дарси определяются при учете силы тяжести по формулам
\begin{equation}\begin{aligned}
w_\text{в}&=-\frac{k_1}{\mu_\text{в}}\left(\frac{\partial p}{\partial s}+\rho_\text{в}g\frac{\partial z}{\partial s}\right);&w_\text{н}&=-\frac{k_2}{\mu_\text{н}}\left(\frac{\partial p}{\partial s}+\rho_\text{н}g\frac{\partial z}{\partial s}\right).
\end{aligned}\tag{8.39}\label{q9:eq-8-39}\end{equation}
Вследствие неизбежных возмущений на границе раздела частицы воды попадают в область, занятую нефтью, при этом их дальнейшее движение может либо ускоряться, либо замедляться. В первом случае, при ускорении движения частиц воды, движение границы раздела будет неустойчивым; во втором, при замедлении движения частиц воды, --- устойчивым.


Условия устойчивости движения границы раздела можно установить из следующих элементарных соображений. Обозначим через $(w_\text{в})_\text{н}$ скорость фильтрации частиц воды, попавших в поток нефти с градиентом давления $(\partial p/\partial s)_\text{н}$; через $(k_1)_2$ --- проницаемость пласта для воды в зоне движения нефти.

Тогда из первого соотношения (8.39) имеем
\begin{equation}(w_\text{в})_\text{н}=-\frac{(k_1)_2}{\mu_\text{в}}\left[\left(\frac{\partial p}{\partial s}\right)_\text{н}+\rho_\text{в}g\frac{\partial z}{\partial s}\right].\tag{8.40}\label{q9:eq-8-40}\end{equation}
Скорость фильтрации основных частиц нефти, соприкасающихся с проникшими туда частицами воды, согласно второму уравнению (8.39) будет
\begin{equation}w_\text{н}=-\frac{k_2}{\mu_\text{н}}\left[\left(\frac{\partial p}{\partial s}\right)_\text{н}+\rho_\text{н}g\frac{\partial z}{\partial s}\right].\tag{8.41}\label{q9:eq-8-41}\end{equation}
Из уравнений (8.40) и (8.41) получаем связь между скоростями фильтрации $(w_\text{в})_\text{н}$ и $w_\text{н}$:
\[
\frac{\mu_\text{в}}{(k_1)_2}(w_\text{в})_\text{н}+\rho_\text{в}g\frac{\partial z}{\partial s}=\frac{\mu_\text{н}}{k_2}w_\text{н}+\rho_\text{н}g\frac{\partial z}{\partial s}=-\left(\frac{\partial p}{\partial s}\right)_\text{н},
\]
откуда
\[
(w_\text{в})_\text{н}=\frac{\mu_\text{н}}{\mu_\text{в}}\frac{(k_1)_2}{k_2}w_\text{н}-\frac{(k_1)_2}{\mu_\text{в}}(\rho_\text{в}-\rho_\text{н})g\frac{\partial z}{\partial s}.
\]
Об устойчивости движения границы раздела можно судить по разности скоростей фильтрации:
\begin{equation}\begin{aligned}
\Delta w&=(w_\text{в})_\text{н}-w_\text{н}=\left[\frac{\mu_\text{н}}{\mu_\text{в}}\frac{(k_1)_2}{k_2}-1\right]w_\text{н}-\\
&\quad-\frac{(k_1)_2}{\mu_\text{в}}(\rho_\text{в}-\rho_\text{н})g\frac{\vert\partial z}{\partial s}.
\end{aligned}\tag{8.42}\label{q9:eq-8-42}\end{equation}
При $\Delta w\le0$ движение границы раздела жидкостей будет устойчивым, при $\Delta w\ge0$ движение неустойчиво.

Если угол наклона пласта к горизонту обозначить через $\alpha$, то, очевидно, $\partial z/\partial s=\sin\alpha$.

Тогда условие устойчивости границы раздела (8.42) можно представить в виде
\[
\Delta w=\left[\frac{\mu_\text{н}}{\mu_\text{в}}\frac{(k_1)_2}{k_2}-1\right]w_\text{н}-\frac{(k_1)_2}{\mu_\text{в}}(\rho_\text{в}-\rho_\text{н})g\sin\alpha.
\]
Обычно $(k_1)_2$ меньше $k_2$. Считая в первом приближении, что $(k_1)_2=k_2$, преобразуем это соотношение к виду
\begin{equation}\Delta w=\left(\frac{\mu_\text{н}}{\mu_\text{в}}-1\right)w_\text{н}-\frac{k_2}{\mu_\text{в}}(\rho_\text{в}-\rho_\text{н})g\sin\alpha.\tag{8.43}\label{q9:eq-8-43}\end{equation}
Так как при устойчивом движении границы раздела $\Delta w\le0$, то из (8.43) найдем, что при устойчивом движении границы раздела жидкостей скорость фильтрации нефти $w_\text{н}$ на границе раздела должна быть
\[
w_\text{н}\le\frac{k_2}{\mu_\text{н}-\mu_\text{в}}(\rho_\text{в}-\rho_\text{н})g\sin\alpha.
\]
Более строгое исследование рассмотренной задачи проводится методами теории возмущений и гидродинамической устойчивости.
\par\nobreak\smallskip
{\small К. С. Басниев, А. М. Власов, И. Н. Кочина, В. М. Максимов. \textit{Подземная гидравлика}. 1986. Глава 8, § 4, с. 197--198.}\par
\end{quote}

\noindent В формуле (8.42) перед $\partial z$ на скане виден дополнительный вертикальный штрих. Это место требует сверки с более чётким изображением оригинала; штрих сохранён.
\end{document}
